\section{衍射积分中的尺度与归一化}
\subsection{几何小量不等于相位小量}
传播距离可展开为$\sqrt{L^2+\rho^2}=L+\rho^2/(2L)-\rho^4/(8L^3)+\cdots$。在幅度因子$1/r$中，只要$\rho/L\ll1$就可近似$r=L$；在相位$kr$中，误差还乘$k=2\pi/\lambda$，要求$k\rho^4/(8L^3)$足够小才能只保留二次项。把三角近轴条件与相位截断条件分开，才能判断菲涅耳近似在特定尺寸下的精度。

夫琅禾费近似还要求孔面二次相位$k(\xi^2+\eta^2)/(2L)$的变化足够小，或已由透镜补偿。增大屏距、缩小孔径都可减小该项；仅让观察角很小还不够。高斯束、直边及大孔径近场计算通常仍须保留二次相位。

在孔径前还有有限距离点光源时，入射到孔面已带一个二次球面相位。传播到屏的二次相位与它相加，等效距离由两段波前曲率共同决定。不能只看孔到屏的$L$来断定远场；若点光源很近，入射曲率项可能仍很大。理想透镜可用反号二次相位补偿相应曲率，形成清楚的傅里叶变换面。

\subsection{远场坐标的量纲}
一维远场场可写$U(x)=C F(q)$，其中$F(q)=\int U_0(\xi)\ee^{-iq\xi}\dd\xi$，$q=kx/L$。$q$量纲为长度倒数，$x$是屏幕长度。若使用空间频率$f_x=x/(\lambda L)$，指数应写$\ee^{-i2\pi f_x\xi}$。这两套约定只差$2\pi$，但单缝零点和高斯宽度的常数必须随之同步。

\Diagram{fig:o3scales}{孔面长度尺度与远场空间频率尺度互为倒数}{\draw[axis](0,1.4)--(5,1.4)node[right]{$\xi$};\draw[very thick](1.5,1.4)--(3.5,1.4);\node[above]at(2.5,1.4){宽$a$};\draw[axis](0,0)--(5,0)node[right]{$q$};\draw[very thick](2,0)--(3,0);\node[below]at(2.5,0){宽度尺度$1/a$};}
图\ref{fig:o3scales}说明孔越宽、角谱越窄，但比较绝对亮度还要保留孔面照明归一化。均匀单缝振幅正比$a$，所以中央强度随$a^2$增长。把每条曲线都除以各自峰值后只剩形状，不能再从归一化图上读取绝对功率变化。

\section{傅里叶积分中的能量守恒}
对平方可积孔场，采用上述无$2\pi$的傅里叶约定，帕塞瓦尔关系为
\[
 \int_{-\infty}^{\infty}|F(q)|^2\dd q
 =2\pi\int_{-\infty}^{\infty}|U_0(\xi)|^2\dd\xi.
\]
它可从双重积分展开$|F|^2$后，使用$\int\ee^{-iq(\xi-\eta)}\dd q=2\pi\delta(\xi-\eta)$得到。自由传播核的系数与屏面变量$\dd x=L\dd q/k$相配合，使积分能流等于孔面透过能流。二维情况沿两个方向各用一次同样关系。

单缝可直接核对功率归一化：孔面振幅恒定时，$F(q)=aU_0\operatorname{sinc}(qa/2)$。换元$u=qa/2$，得$\int|F|^2\dd q=2a|U_0|^2\int\operatorname{sinc}^2u\dd u$，积分正比于$a$。孔宽增大使中央强度升高、角分布变窄，二者共同满足总功率守恒。

\subsection{相位孔与振幅孔}
若孔面仅乘单位模相位$\ee^{i\phi(\xi)}$，总透过功率不变，但远场分布可明显改变。线性相位整体转向，二次相位改变聚焦或观察面，周期相位形成多个衍射级。若孔面乘振幅透过率$|t(\xi)|<1$，则总功率减少；它仍可通过压低边缘权重改善旁瓣，代价是峰宽或通光量。

多束加权的二项式系数和连续高斯照明都是这种思想的离散、连续例子。比较不同孔径时应固定哪一个条件：总输入功率、最大照明强度、有效孔径宽度，还是透过总能量。若约束未说明，峰高、旁瓣与“效率”的比较往往没有唯一答案。

\section{点像、物像与分辨判据}
\subsection{单点响应怎样组合成像}
理想无像差透镜的几何像点在波动理论中仍有有限衍射斑。孔函数的傅里叶变换给出点源的复振幅响应$h(x)$，其强度响应为$|h(x)|^2$。多个彼此不相干物点成像时，各点像强度相加；若物点相干，则应先把复振幅相加。两者即使几何点位完全相同，图样也可能因相对相位不同而改变。

\Diagram{fig:o3pointimages}{不相干物点的点像相加强度，相干物点则先加场}{\fill(0,.7)circle(1.5pt)node[left]{$O_1$};\fill(0,-.7)circle(1.5pt)node[left]{$O_2$};\draw[<->](2,-1.5)--(2,1.5);\draw[curve](4.5,.7)circle(.45);\draw[curve](4.5,-.7)circle(.45);\draw[guide](0,.7)--(4.5,-.7);\draw[guide](0,-.7)--(4.5,.7);}
图\ref{fig:o3pointimages}只标出点像中心与宽度。对两个等强不相干谱线或星点，瑞利判据使用主峰与首暗环重合；对两相干点，交叉项会使中间强度取决于相位，不能沿用同一张强度相加曲线。

\subsection{瑞利、曲率和测量误差}
瑞利判据以主峰与首零点重合作为分离尺度。曲率判据则考察中央由峰变谷的边界，在对称等强模型中令二阶导数为零。对sinc平方线形，后者间距约为瑞利尺度的.8296倍；高斯、圆孔等线形须另求这个系数。

实际谱线很弱时，另一个强峰的翼部与背景噪声可能比中心曲率更重要。高信噪比、已知线形的拟合有时能估计小于传统瑞利间隔的峰位，但这不等于光学点像变窄；估计不确定度取决于光子数、噪声、模型先验和参数相关。竞赛题若指定某种判据，应按该定义计算，不把不同定义混作唯一“最小可分辨距离”。

\subsection{光栅参数中的三个长度}
缝宽$a$、周期$b$、照明总宽$W\simeq Nb$决定不同结构。$a$给单缝包络和缺级；$b$给各级中心角和最高可传播级次；$W$给同级窄峰的宽度。若固定$W$而缩小$b$，虽然$N$增大，可用级次却下降；若固定$b$而增加$N$，主峰中心不动，宽度变窄。

闪耀小面的倾角把单元包络峰移向选定级次，是效率设计；增加总照明宽度才真正增加可参与的相干横向范围。真实光栅还受表面误差、局部周期误差和照明相位不均影响，实际窄峰可能比理想$mN$更宽。这些误差应按具体分布进入孔面场，再求实际线形。

\subsection{孔径、相干照明和点像强度}
夫琅禾费场是孔面复场的傅里叶变换，而非孔面强度的傅里叶变换。若孔面透过率为$t(\xi)$、照明复场为$A(\xi)$，实际参与积分的是$A(\xi)t(\xi)$。均匀照明、振幅加权和附加相位因而改变不同对象：加权改变面元振幅，相位片改变面元方向，遮挡则同时改变通光范围。

对若干点源成像，还要在像面区分相干与不相干叠加。若两个点源经同一系统形成复振幅点响应$h(x-x_1)$、$h(x-x_2)$，完全相干时强度为
\[
 |a_1h(x-x_1)+a_2\ee^{i\phi}h(x-x_2)|^2.
\]
互不相干时交叉项长期平均为零，才得到$|a_1|^2|h(x-x_1)|^2+|a_2|^2|h(x-x_2)|^2$。瑞利判据和本讲等强双峰曲率判据采用后一种模型；在孔径和点源间距固定时，相干性仍会改变中央峰谷。
\Diagram{fig:o3coherentpairs}{实线为同相合场，曲线虚线为反相合场；竖直虚线标出点像中心}{\draw[axis](-.3,0)--(5.7,0)node[right]{$x$};\draw[axis](0,-1)--(0,2.7)node[above]{场};\draw[guide](2,-.4)--(2,2.3);\draw[guide](3.2,-.4)--(3.2,2.3);\draw[curve,domain=.2:5.2,samples=100]plot(\x,{exp(-2*(\x-2)^2)+exp(-2*(\x-3.2)^2)});\draw[dashed,domain=.2:5.2,samples=100]plot(\x,{exp(-2*(\x-2)^2)-exp(-2*(\x-3.2)^2)});\node at(4.7,2){同相};\node at(4.7,-.7){反相};}
图\ref{fig:o3coherentpairs}为平滑实场的示意，实线同相相加，虚线反相相加；它不代替具体圆孔或单缝积分。使用测得的强度曲线讨论分辨时，应先明确各点响应、光源统计和探测平均方式，再选择判据。

远场强度还能检验算式归一化：对完整角域或相应近轴傅里叶平面积分，应与孔面透过功率相符。若只截取主瓣，则所得只是该区域功率；忽略旁瓣后不能继续把主瓣积分称为总功率。若孔径尺寸接近波长、观察角接近掠射或介质边界引入明显矢量效应，标量近轴模型的适用性也须重新判断。
